\section{Marco Te\'orico}

\subsection{Computaci\'on Cu\'antica}

\subsubsection{Qubits y estados cu\'anticos}
Un qubit es la unidad fundamental de informaci\'on cu\'antica. A diferencia del bit cl\'asico que toma valores 0 o 1, un qubit puede existir en una superposici\'on de ambos estados:
\begin{equation}
|\psi\rangle = \alpha|0\rangle + \beta|1\rangle, \quad |\alpha|^2 + |\beta|^2 = 1
\end{equation}
donde $\alpha, \beta \in \mathbb{C}$ son amplitudes de probabilidad.

\subsubsection{Puertas cu\'anticas}
Las puertas cu\'anticas son operaciones unitarias que transforman el estado de los qubits. Las puertas de rotaci\'on $R_Y(\theta)$ y $R_Z(\theta)$ son particularmente relevantes para la codificaci\'on de datos:
\begin{equation}
R_Y(\theta) = \begin{bmatrix} \cos\frac{\theta}{2} & -\sin\frac{\theta}{2} \\ \sin\frac{\theta}{2} & \cos\frac{\theta}{2} \end{bmatrix}
\end{equation}
Las puertas CNOT generan entrelazamiento entre qubits, esencial para capturar correlaciones en los datos.

\subsubsection{Circuitos cu\'anticos parametrizados}
Un circuito parametrizado $U(\mathbf{x}, \boldsymbol{\theta})$ consiste en un \textit{feature map} $U_\phi(\mathbf{x})$ que codifica datos cl\'asicos en estados cu\'anticos, y un \textit{ansatz} $U_W(\boldsymbol{\theta})$ con par\'ametros entrenables.

\subsection{Machine Learning Cl\'asico}

\subsubsection{Support Vector Machine (SVM)}
SVM busca el hiperplano que maximiza el margen entre clases. Con el kernel RBF $K(\mathbf{x}_i, \mathbf{x}_j) = \exp(-\gamma\|\mathbf{x}_i - \mathbf{x}_j\|^2)$, SVM clasifica datos no linealmente separables.

\subsubsection{MLP y KNN}
El Perceptr\'on Multicapa (MLP) aprende funciones no lineales mediante capas ocultas. KNN clasifica cada muestra seg\'un la clase mayoritaria de sus $k$ vecinos m\'as cercanos.

\subsection{Quantum Machine Learning}

\subsubsection{Variational Quantum Classifier (VQC)}
El VQC \citep{havlicek2019} utiliza un circuito parametrizado con entrenamiento h\'ibrido cu\'antico-cl\'asico. La funci\'on de costo es la entrop\'ia cruzada:
\begin{equation}
\mathcal{L}(\boldsymbol{\theta}) = -\frac{1}{N}\sum_{i=1}^{N}\sum_{c=1}^{C} y_{ic} \log(\hat{y}_{ic}(\boldsymbol{\theta}))
\end{equation}

\subsubsection{Quantum SVM (QSVC)}
QSVC \citep{rebentrost2014} reemplaza el kernel cl\'asico por un kernel cu\'antico basado en la fidelidad entre estados:
\begin{equation}
K_Q(\mathbf{x}_i, \mathbf{x}_j) = |\langle\phi(\mathbf{x}_j)|\phi(\mathbf{x}_i)\rangle|^2
\end{equation}

\subsubsection{Feature Maps}
Los feature maps codifican datos cl\'asicos en el espacio cu\'antico:
\begin{itemize}
    \item \textbf{ZFeatureMap}: Rotaciones $R_Z(x_i)$ sin entrelazamiento.
    \item \textbf{ZZFeatureMap}: Interacciones $R_{ZZ}(x_i \cdot x_j)$, capturando correlaciones de segundo orden.
    \item \textbf{PauliFeatureMap}: Generaliza los anteriores con cadenas de Pauli.
\end{itemize}

\subsubsection{Ansatz y Barren Plateaus}
El ansatz RealAmplitudes consiste en capas de rotaciones $R_Y$ y puertas CNOT. El par\'ametro \textit{reps} controla la profundidad. Circuitos demasiado profundos pueden sufrir de \textit{barren plateaus}: regiones donde los gradientes se desvanecen exponencialmente \citep{cerezo2021}.

\subsubsection{Optimizadores}
Los optimizadores cl\'asicos m\'as utilizados en VQC son:
\begin{itemize}
    \item \textbf{COBYLA}: M\'etodo libre de gradiente, estable para simulaciones.
    \item \textbf{SPSA}: Aproxima gradientes con perturbaciones estoc\'asticas, robusto al ruido.
    \item \textbf{L-BFGS-B}: M\'etodo quasi-Newton, requiere gradientes exactos.
\end{itemize}
